\documentclass{article}
\usepackage[a4paper, total={7in, 9.2in}]{geometry}
\usepackage{titling}
\usepackage{tabularx}
\usepackage{hyperref}
\usepackage{array}
\usepackage{longtable}
\usepackage[table]{xcolor}
\usepackage{booktabs}
\usepackage{float}
\usepackage{fancyhdr}

\pagestyle{fancy}
\fancyhf{}
\renewcommand{\headrulewidth}{0.4pt}
\fancyhead[C]{\textsc{The National University of Computer and Emerging Sciences}}
\fancyfoot[C]{\thepage}

\title{
{\Large\scshape The National University of Computer and Emerging Sciences}\\[10pt]
{\large Course Outline}\\[4pt]
{\Huge\textbf{Data Compression CS-4084}}\\[8pt]
{\normalsize Semester Spring-2026}
}
\author{ }
\date{}

\pretitle{\begin{center}}
\posttitle{\end{center}}
\preauthor{}
\postauthor{}
\predate{}
\postdate{}

\begin{document}

\maketitle
\vspace{-0.5cm}

\begin{table}[H]
    \centering
    \def\arraystretch{1.5}
    \begin{tabularx}{\textwidth}{ l X l X }
        \textbf{Instructor}: & Dr. Faisal Aslam & \textbf{Office hours}: & Mon/Wed 14:00 -- 15:00 \\
        \textbf{Credit Hours:} & 3 & \textbf{Prerequisites:} & Data Structures \& Algorithms\\
        & & & Probability \& Discrete Mathematics \\
    \end{tabularx}
    \label{tab:Intro}
\end{table}
\vspace{-1cm}

\section*{Course Description}
This course explores the theoretical foundations and algorithmic techniques of data compression, combining principles from information theory, statistical modeling, and practical compressor design. The course progresses from Shannon's entropy-based models through statistical and adaptive entropy coders, dictionary-based methods, redundancy-reduction techniques, transform-based methods, and finally transform coding as applied to lossy image compression (JPEG). Emphasis is placed on understanding both the mathematical limits of compression and the engineering trade-offs in real-world compressors such as DEFLATE, bzip2, and Zstandard, and in widely deployed formats such as JPEG. By the end of the course, students should be able to determine theoretical compression limits for a given source, implement and compare algorithms from different compression paradigms, and evaluate real-world compression systems against application requirements.

\section*{Program Learning Outcomes (PLOs)}
This course covers the following PLOs:

\renewcommand{\arraystretch}{1.4}
\begin{longtable}{|l|p{3.3cm}|p{10.3cm}|}
\hline
\textbf{PLO\#}  & \textbf{PLO Name} & \textbf{PLO Description} \\ \hline
\endfirsthead
\hline
\textbf{PLO\#}  & \textbf{PLO Name} & \textbf{PLO Description} \\ \hline
\endhead
PLO 2 & Knowledge for Solving Computing Problems &
Apply knowledge of computing fundamentals, knowledge of a computing specialization, and mathematics, science, and domain knowledge appropriate for the computing specialization to the abstraction and conceptualization of computing models from defined problems and requirements \\ \hline
PLO 3 & Problem Analysis & Identify, formulate, research literature, and solve complex computing problems reaching substantiated conclusions using fundamental principles of mathematics, computing sciences, and relevant domain disciplines \\ \hline
PLO 4 & Design/Development of Solutions & Design and evaluate solutions for complex computing problems, and design and evaluate systems, components, or processes that meet specified needs with appropriate consideration for public health and safety, cultural, societal, and environmental considerations \\ \hline
\end{longtable}

\section*{Course Learning Outcomes (CLOs)}
CS-4084 is an elective Computer Science course. The course learning outcomes are (university-approved; unchanged):

\renewcommand{\arraystretch}{1.4}
\begin{longtable}{|l|p{9cm}|p{2.7cm}|p{2.2cm}|}
\hline
\textbf{CLO \#}  & \textbf{CLO Description} & \textbf{BT Level} & \textbf{PLO \#}\\ \hline
\endfirsthead
\hline
\textbf{CLO \#}  & \textbf{CLO Description} & \textbf{BT Level} & \textbf{PLO \#}\\ \hline
\endhead
CLO 1 & Analyze and apply fundamental compression theory to data sources. Utilize mathematical frameworks (information theory, rate-distortion theory, algorithmic complexity) to evaluate data redundancy, determine theoretical compression limits, and justify coding scheme selection for specific applications. & Analyze (BT Level 4) & PLO 2, PLO 3 \\ \hline
CLO 2 & Implement and compare algorithmic approaches across compression paradigms. Design, code, and experiment with core compression algorithms from different classes (entropy coding, dictionary-based, transform-based, predictive) to understand their mechanisms, implementation challenges, and performance characteristics. & Create (BT Level 6) & PLO 2, PLO 4 \\ \hline
CLO 3 & Evaluate real-world compressors using appropriate metrics and methodologies. Benchmark and analyze compression systems considering multiple performance dimensions (compression ratio, speed, memory, distortion where applicable), and recommend suitable solutions for given application constraints and data types. & Evaluate (BT Level 5) & PLO 3, PLO 4 \\ \hline
CLO 4 & Assess advanced compression techniques and emerging trends in the field. Critique sophisticated modeling methods and emerging approaches, discussing their theoretical foundations, practical implementations, advantages, limitations, and potential evolution within the compression landscape. & Evaluate (BT Level 5) & PLO 2, PLO 3 \\ \hline
\end{longtable}

\section*{Primary Textbooks}
These are the core references for lecture preparation and are cited by short code throughout the lecture table below.
\begin{enumerate}
    \item[{[IDC]}] Sayood, K. (2017). \textit{Introduction to Data Compression} (5th ed.). Morgan Kaufmann. --- primary reference for entropy coding, dictionary methods, and general compression theory.
    \item[{[EIT]}] Cover, T. M., \& Thomas, J. A. (2006). \textit{Elements of Information Theory} (2nd ed.). Wiley. --- primary reference for information-theoretic foundations.
\end{enumerate}

\section*{Supplementary / Further Reading}
A curated set of original papers and standards, selected to give the deepest single treatment of each major paradigm covered in the course; useful for deeper treatment but not required for lecture preparation.
\begin{enumerate}
    \item[{[Sha48]}] Shannon, C. E. (1948). A Mathematical Theory of Communication. \textit{Bell System Technical Journal}, 27(3), 379--423.
    \item[{[WNC87]}] Witten, I. H., Neal, R. M., \& Cleary, J. G. (1987). Arithmetic Coding for Data Compression. \textit{Communications of the ACM}, 30(6), 520--540.
    \item[{[LZ77]}] Ziv, J., \& Lempel, A. (1977). A Universal Algorithm for Sequential Data Compression. \textit{IEEE Transactions on Information Theory}, 23(3), 337--343.
    \item[{[RFC1951]}] Deutsch, P. (1996). \textit{DEFLATE Compressed Data Format Specification version 1.3}. RFC 1951.
    \item[{[BW94]}] Burrows, M., \& Wheeler, D. J. (1994). A Block-Sorting Lossless Data Compression Algorithm. Digital Equipment Corporation Technical Report.
    \item[{[Wal91]}] Wallace, G. K. (1991). The JPEG Still Picture Compression Standard. \textit{IEEE Transactions on Consumer Electronics}, 38(1), xviii--xxxiv.
\end{enumerate}

\section*{Lecture-Wise Course Outline (30 Lectures)}
\renewcommand{\arraystretch}{1.35}
\setlength{\tabcolsep}{4pt}
\footnotesize
\begin{longtable}{|c|p{8.7cm}|p{1.6cm}|p{2.6cm}|}
\hline
\textbf{L\#} & \textbf{Topic} & \textbf{CLO(s)} & \textbf{Reading} \\ \hline
\endfirsthead
\hline
\textbf{L\#} & \textbf{Topic} & \textbf{CLO(s)} & \textbf{Reading} \\ \hline
\endhead

\multicolumn{4}{|c|}{\cellcolor{gray!20}\textbf{Part I: Introduction and Theoretical Foundations (Lectures 1--6)}} \\ \hline
1 & Motivation for data compression; applications; lossless vs.\ lossy compression; benefits of compression & CLO1 & [IDC], [EIT] \\ \hline
2 & Compression performance metrics (size- and rate-based); worked example (audio compression); introduction to Huffman coding & CLO1, CLO2 & [IDC] \\ \hline
3 & Types of codes (nonsingular, uniquely decodable, instantaneous/prefix); basic terminology and notation & CLO1 & [EIT] \\ \hline
4 & Information and redundancy; entropy: definition, calculation, entropy of English text, entropy rate & CLO1 & [Sha48], [EIT] \\ \hline
5 & Entropy as a lower bound; the Kraft--McMillan inequality; optimality of Huffman codes (theory) & CLO1 & [EIT] \\ \hline
6 & Limitations of symbol-by-symbol coding; block coding and its trade-offs; bridging to arithmetic coding and context modeling & CLO1 & [EIT], [IDC] \\ \hline

\multicolumn{4}{|c|}{\cellcolor{gray!20}\textbf{Part II: Advanced Entropy Coding (Lectures 7--9)}} \\ \hline
7 & Coding taxonomy and framework; Shannon coding (1948) & CLO1, CLO2 & [Sha48], [IDC] \\ \hline
8 & Shannon--Fano coding (1949); Canonical Huffman codes & CLO2 & [IDC] \\ \hline
9 & Adaptive Huffman coding (FGK, Vitter); Arithmetic coding: the paradigm shift; comparison and synthesis of entropy coders & CLO2, CLO3 & [IDC], [WNC87] \\ \hline

\multicolumn{4}{|c|}{\cellcolor{gray!20}\textbf{Part III: Source Modeling and Statistical Dependence (Lectures 10--11)}} \\ \hline
10 & Memoryless sources vs.\ sources with memory; conditional entropy and mutual information; Markov sources & CLO1 & [EIT] \\ \hline
\multicolumn{4}{|c|}{\cellcolor{yellow!25}\textit{--- Midterm 1 (after Lecture 10; covers L1--L10) ---}} \\ \hline
11 & Entropy rate revisited; context modeling in practice; case study: text compression modeling; the modeling--coding separation principle; adaptive vs.\ static modeling & CLO1, CLO3 & [IDC] \\ \hline

\multicolumn{4}{|c|}{\cellcolor{gray!20}\textbf{Part IV: Dictionary-Based Compression and DEFLATE (Lectures 12--17)}} \\ \hline
12 & Limits of statistical coding; the paradigm shift to dictionary methods; LZ77: the sliding-window algorithm & CLO2 & [LZ77] \\ \hline
13 & LZSS: practical improvements to LZ77; LZ78: the dictionary-growth algorithm; LZW and its applications (GIF, UNIX compress) & CLO2 & [IDC], [LZ77] \\ \hline
14 & Theoretical properties of Lempel--Ziv methods: universality and asymptotic optimality; bridging dictionary and entropy coding (hybrid approaches) & CLO1 & [LZ77], [EIT] \\ \hline
15 & DEFLATE: history and architecture (two-stage pipeline); LZSS in DEFLATE (32KB window, hash chains, lazy matching) & CLO2 & [RFC1951] \\ \hline
16 & DEFLATE's Huffman coding scheme (literal/distance alphabets, canonical Huffman via code lengths); block structure and streaming format & CLO2 & [RFC1951] \\ \hline
17 & Step-by-step DEFLATE compression walkthrough; compression levels and trade-offs; applications (gzip, ZIP, PNG, HTTP); limitations; modern alternatives (LZMA, Zstandard, Brotli, LZ4) & CLO2, CLO3, CLO4 & [RFC1951], [IDC] \\ \hline

\multicolumn{4}{|c|}{\cellcolor{gray!20}\textbf{Part V: Redundancy Reduction and Transform-Based Compression (Lectures 18--23)}} \\ \hline
18 & Run-length encoding (RLE): fundamental concepts and theory of runs; taxonomy of RLE methods; count-based RLE (fixed- and variable-width) & CLO2 & [IDC] \\ \hline
19 & Escape-based RLE; block/header-based RLE (PackBits); zero RLE (JPEG case study); comparison of variants; RLE as a preprocessing technique & CLO2, CLO3 & [IDC] \\ \hline
20 & The three paradigms of lossless compression; motivation for global reordering; the Burrows--Wheeler Transform (intuition and step-by-step example) & CLO1 & [BW94] \\ \hline
\multicolumn{4}{|c|}{\cellcolor{yellow!25}\textit{--- Midterm 2 (after Lecture 20; covers L11--L20) ---}} \\ \hline
21 & Formal definition of the BWT; inverting the BWT (LF-mapping, reconstruction algorithm, correctness) & CLO1, CLO2 & [BW94] \\ \hline
22 & Why the BWT improves compressibility; the full BWT pipeline (BWT + MTF + RLE + entropy coding); efficient implementation via suffix arrays & CLO2 & [BW94] \\ \hline
23 & Comparison with Lempel--Ziv methods; case study: the bzip2 pipeline; limitations and modern extensions (e.g., BSC) & CLO3, CLO4 & [BW94], [IDC] \\ \hline

\multicolumn{4}{|c|}{\cellcolor{gray!20}\textbf{Part VI: Advanced Statistical and Entropy Coding (Lectures 24--29)}} \\ \hline
24 & Prediction by Partial Matching (PPM): beyond finite-context models; the PPM algorithm; the escape and prediction mechanism; PPMC escape estimation; the exclusion principle & CLO2 & [IDC] \\ \hline
25 & Step-by-step PPM examples; implementation challenges (memory management, context tries, the zero-frequency problem) & CLO2 & [IDC] \\ \hline
26 & PPM variants and successors (PPM*, PPMZ, PAQ context mixing); why PPM matters today & CLO4 & [IDC] \\ \hline
27 & Arithmetic coding: theory, algorithm, and a step-by-step worked example & CLO2 & [WNC87] \\ \hline
28 & Range coding: theory, algorithm, renormalisation, and decoding; arithmetic vs.\ range coding comparison & CLO2 & [WNC87], [IDC] \\ \hline
29 & Asymmetric Numeral Systems (ANS): motivation and central idea; rANS encoding/decoding; tANS (encode/decode tables); ANS vs.\ Arithmetic Coding vs.\ Huffman; ANS in practice (Zstandard) & CLO2, CLO3, CLO4 & [IDC] \\ \hline

\multicolumn{4}{|c|}{\cellcolor{gray!20}\textbf{Part VII: Transform Coding for Images (Lecture 30)}} \\ \hline
30 & Transform coding motivation; the Discrete Cosine Transform; the complete JPEG pipeline (level shift, 2D DCT, quantization, zig-zag scan, RLE, Huffman coding); rate--distortion trade-off; compression artefacts; JPEG vs.\ modern formats & CLO2, CLO3 & [Wal91] \\ \hline
\multicolumn{4}{|c|}{\cellcolor{yellow!25}\textit{--- Comprehensive Final Exam (covers L1--L30) ---}} \\ \hline
\end{longtable}

\normalsize

\section*{Assessment Instruments}
\begin{itemize}
    \item \textbf{Assignments} -- problem sets and short implementation tasks distributed across the semester, covering entropy coding, dictionary-based methods, and transform-based techniques.
    \item \textbf{Quizzes} -- short quizzes distributed across the semester, primarily on the lecture(s) immediately preceding each quiz.
    \item \textbf{Midterm 1} (after L10), \textbf{Midterm 2} (after L20), \textbf{Comprehensive Final} (after L30).
\end{itemize}

\section*{Grading Scheme}
\begin{itemize}
    \item Assignments -- 5\%
    \item Quizzes -- 15\%
    \item Midterm-1 -- 15\%
    \item Midterm-2 -- 15\%
    \item Final Exam -- 50\%
\end{itemize}

\end{document}
